\chapter{Family--defect represented degree-two junk gate}

Предыдущий audit свёл classical parent problem к точному вопросу: является
ли
\begin{equation}
 \mu_G=XX^T-|\Phi|^2I_3
\end{equation}
канонической self-adjoint two-form explicit 18-dimensional KO6 geometry.
Здесь вычисляется не blockwise analogy, а полный represented universal
calculus
\begin{equation}
 \Omega_D^2(\mathcal A_{\rm q})
 =\pi(\Omega_u^2\mathcal A_{\rm q})/\pi(d\ker\pi_1)
\end{equation}
для
\(\mathcal A_{\rm q}=\mathbb R_0\oplus M_3(\mathbb R)_G\oplus\mathbb C_2\).
Именно quotient по differential ideal является стандартным определением
Connes forms; см. \path{arXiv:hep-th/9605001}. Конечномерные three-point
calculi и явный quotient также разобраны в \path{arXiv:2405.07866}.

\section{Алгоритм quotient}

Для real basis \(e_a\) алгебры строятся generators
\begin{align}
 \omega_{ab}&=\pi(e_a)[D_F,\pi(e_b)],\\
 \eta_{abc}&=\pi(e_a)[D_F,\pi(e_b)][D_F,\pi(e_c)].
\end{align}
Если coefficients \(c_{ab}\) удовлетворяют
\begin{equation}
 \sum_{ab}c_{ab}\omega_{ab}=0,
\end{equation}
то corresponding degree-two junk равен
\begin{equation}
 \sum_{ab}c_{ab}[D_F,\pi(e_a)][D_F,\pi(e_b)].
\end{equation}
Все spans, kernels и quotient ranks вычисляются SVD. Реализация сначала
проверена на existing \(\mathbb C^3\) path control: length-two endpoint
matrix units снова оказываются junk с residual меньше \(10^{-15}\).

\section{Generic KO6 ledger}

На четырёх независимых full-rank real matrices \(X\) при nonzero complex
\(\Phi\) complexified calculus стабильно даёт
\begin{equation}
 \rank\pi(\Omega_u^1)=20,
 \qquad
 \rank\pi(\Omega_u^2)=20,
 \qquad
 \rank\pi(d\ker\pi_1)=9,
\end{equation}
и потому
\begin{equation}
 \rank\Omega_D^2=11.
 \label{eq:family-degree-two-quotient-rank}
\end{equation}
Junk лежит внутри represented two-forms с residual меньше
\(3\times10^{-15}\). Независимый real-linear audit даёт ranks
\begin{equation}
 22,
 \qquad 22,
 \qquad 9,
 \qquad 13,
\end{equation}
что является realification той же obstruction, а не восстановлением
six-dimensional symmetric sector.

\section{Что происходит с middle curvature}

На particle half represented middle diagonal block имеет полный rank девять,
но его junk image также имеет rank девять. Следовательно,
\begin{equation}
 \boxed{
 \operatorname{pr}_{G,\rm particle}\Omega_D^2=0.}
 \label{eq:family-particle-middle-twoform-zero}
\end{equation}
В полном KO6 doubling сохраняется одна complex middle direction, однако она
находится только на conjugate chain и пропорциональна \(I_3\):
\begin{equation}
 \rank_{\mathbb C}\operatorname{pr}_{G,\rm conj}\Omega_D^2=1,
 \qquad
 \rank_{\mathbb C}operatorname{pr}_{\operatorname{Sym}_{3,0}}
 \Omega_D^2=0.
 \label{eq:family-only-central-conjugate-twoform}
\end{equation}
Таким образом, five-dimensional traceless shape curvature исчезает, а
полный module
\(\operatorname{Sym}_3(\mathbb R)=\mathbf1\oplus\mathbf5\) не возникает.
Оставшаяся central conjugate direction не может воспроизвести
\(XX^T-|\Phi|^2I_3\) и не имеет требуемой particle/conjugate symmetry.

\begin{nogocriterion}[Ordinary degree-two symmetric-curvature no-go]
Для explicit KO6 representation обычный represented two-form quotient не
содержит six-dimensional middle self-adjoint curvature. На physical
particle chain весь middle \(M_3\) удаляется junk ideal; в doubled quotient
остаётся только одна central conjugate direction. Поэтому
\(\tau_3(\mu_G^2)\) нельзя вывести как norm ordinary Connes two-form этой
geometry.
\end{nogocriterion}

\section{Почему radial rank jump не спасает route}

При специальном background \(X=\rho I_3\) commutators вырождаются и
degree-two junk rank падает до нуля. Тогда middle matrices формально снова
видны. Но differential calculus parent theory должен быть определён
off-shell и устойчив на открытом множестве configurations. Module,
возникающий только после подстановки уже выбранного isotropic vacuum, не
может служить derivation самого shape-locking potential. Поэтому radial
rank jump является singular-locus artifact, а не rescue.

\section{Оставшийся путь}

После этого gate classical цепочка исчерпана:
\begin{equation}
 \text{ordinary trace}
 \;\longrightarrow\;\text{wrong sign},
 \qquad
 \text{mapping cone}
 \;\longrightarrow\;\text{endpoint only},
\end{equation}
\begin{equation}
 \text{real auxiliary}
 \;\longrightarrow\;\text{wrong sign},
 \qquad
 \text{ordinary }\Omega_D^2/J
 \;\longrightarrow\;\text{middle shape is junk}.
\end{equation}
Условно открыта только quantum-measure identity с imaginary
Hubbard--Stratonovich coupling. Литература подтверждает, что repulsive
channels требуют imaginary coupling и complex saddle; см.
\path{arXiv:2308.05150}. Для project pass теперь необходимо вывести этот
contour и determinant/Pfaffian weight из той же KO6 fermionic measure.
Иначе потребуется modified differential calculus, twist или новый finite
carrier, что уже является новой architecture, а не продолжением current
ordinary spectral triple.

Машинный ledger сохранён в
\path{s2t_v4_family_defect_degree_two_junk_gate_results.json}; генератор ---
\path{s2t_v4_family_defect_degree_two_junk_gate.py}.